\section{Tarefa 5.3: Projeto e montagem de um subtrator completo de 1 bit}

Nesta etapa, foi solicitado projetar e montar um subtrator completo de 1 bit com propagacao do empresta-um. A celula basica deve produzir a diferenca $S_i$ e o sinal de emprestimo $E_{i-1}$ a partir das entradas $A_i$, $B_i$ e $E_i$, utilizando somente portas NAND de duas entradas na implementacao final.

\subsection{Tabela-verdade}

\begin{table}[H]
    \centering
    \caption{Tabela-verdade da celula do subtrator completo.}
    \begin{tabular}{ccccc}
        \toprule
        $A_i$ & $B_i$ & $E_i$ & $S_i$ & $E_{i-1}$ \\
        \midrule
        0 & 0 & 0 & & \\
        0 & 0 & 1 & & \\
        0 & 1 & 0 & & \\
        0 & 1 & 1 & & \\
        1 & 0 & 0 & & \\
        1 & 0 & 1 & & \\
        1 & 1 & 0 & & \\
        1 & 1 & 1 & & \\
        \bottomrule
    \end{tabular}
\end{table}

\subsection{Expressoes logicas}

\[
S_i = \underline{\hspace{8cm}}
\]

\[
E_{i-1} = \underline{\hspace{8cm}}
\]

\subsection{Implementacao usando NAND2}

Mostre aqui a minimizacao das expressoes e a reescrita em uma forma adequada para implementacao apenas com portas NAND de duas entradas. Se desejar, inclua mapas de Karnaugh ou transformacoes algebricas.

\begin{figure}[H]
    \centering
    % \includegraphics[width=0.8\textwidth]{questao3/subtrator_nand.png}
    \fbox{\rule{0pt}{5cm}\rule{0.8\textwidth}{0pt}}
    \caption{Espaco reservado para o circuito em NAND2 do subtrator completo.}
\end{figure}

\subsection{Resultados experimentais}

\begin{table}[H]
    \centering
    \caption{Resultados experimentais do subtrator completo de 1 bit.}
    \begin{tabular}{ccccc}
        \toprule
        $A_i$ & $B_i$ & $E_i$ & $S_i$ & $E_{i-1}$ \\
        \midrule
        0 & 0 & 0 & & \\
        0 & 0 & 1 & & \\
        0 & 1 & 0 & & \\
        0 & 1 & 1 & & \\
        1 & 0 & 0 & & \\
        1 & 0 & 1 & & \\
        1 & 1 & 0 & & \\
        1 & 1 & 1 & & \\
        \bottomrule
    \end{tabular}
\end{table}

Compare aqui os resultados medidos com os valores esperados teoricamente.
